\documentclass{NoteThesis/note-thesis}

\usepackage[version=4]{mhchem}
\usepackage{chemfig}
\setchemfig{atom sep=2em}

\course{Composite Materials}

\coursetitle{QXU5030 Composite Materials}
\semester{2025B}
\teacher{Prof. Shafique Ahmed}

\begin{document}

\maketitle

\tableofcontents

\newpage

\section{Introduction to Module}

\subsection*{Assessment}

\begin{table}[h]
    \centering
    \begin{tabular}{cc}
        \hline
        Activity & Score\\
        \hline
        Coursework(labs) & 40\%\\
        Final Exam & 60\%\\ 
        \hline
    \end{tabular}
\end{table}

Week 2N+1: Kahoot! quiz with immediate flash feedback provided after each question.

Week 2N: Quiz delivered via QM+ (or paper-based assignment), with structured written feedback provided within one week.

\subsection{Basic Concepts}

\begin{mynote_env}[frametitle=Anisotropy]
    A property of a material show different \textbf{mechanical or physical properties in different directions}.
\end{mynote_env}

\subsubsection*{Man-Made Composite}

\textbf{Polymer Matrix Composites (PMC's)} These are the most common and will be the main
area of discussion in this course. Also known as \textbf{FRP} - Fibre Reinforced Polymers (or Plastics) - these materials use a polymer-based resin as the matrix, and a variety of fibres such as glass, carbon and aramid as the reinforcement.

\textbf{Metal Matrix Composites (MMC's)} - Increasingly found in the automotive industry, these materials use a metal such as aluminium as the matrix, and reinforce it with fibres such as silicon carbide.

\textbf{Ceramic Matrix Composites (CMC's)} - Used in very high temperature environments, these materials use a ceramic as the matrix and reinforce it with short fibres, or whiskers such as those made from silicon carbide and boron nitride.

\subsubsection*{Type of Composites}

\fastfigure{composite_trend.png}

\subsection{Properties of Composites}

\subsubsection*{Factors of Composites Properties}

The properties of the composite are determined by:

\begin{enumerate}
    \item The properties of the \textbf{ fibre}
    \item The properties of the \textbf{resin}
    \item The ratio of fibre to resin (\textbf{Fibre Volume Fraction})
    \item The \textbf{geometry and orientation} of the fibres
\end{enumerate}

\subsubsection*{Specific Strength}

$$
\textbf{Strength-to-weight ratio} = \verb|tensile strength| / \verb|density|
$$

Indicates how strong a material is relative to its weight.

\subsubsection*{Lightness, or $E/\rho ^3$ and Lightweight Structural Design}

$E$ = \textbf{Young's modulus} (stiffness)\\
$E/\rho ^3$ = key parameter for \textbf{bending stiffness} and \textbf{buckling resistance} in weight-sensitive structures

\subsubsection*{Synergism, the "Composite Effect"}

Ideally, \textbf{rule of mixture (RoM)} for a physical property P:

$$
P_{\verb|RoM|} = P_1 V_1 + P_2 V_2\; (V_1+V_2 = 1)
$$

However, \textbf{synergism} can make $P_{\verb|observed|} >> P_{\verb|RoM|}$

And synergism can modify the \textbf{mode of fraction} by weakrening of fibre/matrix interface to prevent \textbf{brittle fraction}.

\fastfigure{fraction_composite.png}

\subsubsection*{Safety}

Composites can absorb more energy per kilo than metals in crash.

\section{Fibre Reinforcements}

\subsection{Introduction}

Most fibre have higher mechanical properties than \textbf{unreinforced resin or matrix}.

The four main factors that govern the fibre's contribution are:

\begin{enumerate}
    \item The basic \textbf{mechanical properties} of the fibre itself.
    \item The \textbf{surface interaction} of fibre and resin (the '\textbf{interface}').
    \item The amount of fibre in the composite ('\textbf{Fibre Volume Fraction}').
    \item The \textbf{orientation} of the fibres in the composite.
\end{enumerate}

\subsubsection*{Important Synthetic Fibres}

\textbf{Carbon} - the elite fibre : stiff(higher than steel), strong, epoxies, good temperature resistance, chemically inert.

\textbf{Aramid} - very strong, but limited to in stiffness relative to carbon. 

\textbf{Glass} - strong, similar to aluminium in stiffness, relatively cheap.

\subsection{Glass Fibre}

\subsubsection*{Glass Fibre Production}

Blending quarry products (sand, kaolin, limestone, colemanite) at $1600^\circ \text{C}$ to form liquid glass. 

The liquid is passed through \textbf{micro-fine bushings} and simultaneously cooled to produce glass fibre filaments from \textbf{8-15 microns} in diameter. 

The filaments are drawn at high speed (up to $4000\, m/min$) together into a \textbf{strand} or \textbf{roving}.

Then, \textbf{sizing} and form yarns.

\subsubsection*{Glass Fibre Types}

\textbf{E-glass (electrical)} - lower alkali content and stronger than A glass (alkali). Good tensile and compressive strength and stiffness, good electrical properties and relatively low cost, but impact resistance relatively poor.

\textbf{C-glass (chemical)} - best resistance to chemical attack. Mainly used in the form of surface tissue in the outer layer of laminates used in chemical and water pipes and tanks.

\textbf{R or S-glass} - higher tensile strength and modulus than E-glass, with better wet strength retention. Developed for aerospace and defence industries, and used in some hard ballistic
armour applications. 

\begin{mynote_env}[frametitle=E-Glass Fibre Forms]
    \textbf{Strand}. A compactly bundle of filaments, less used.\\
    \textbf{Yarns}. A twisted bundle of filaments, widely used.\\
    \textbf{Roving}. A untwisted, loosely associated bundle.
\end{mynote_env}

\subsubsection*{Binders and Sizing}

The sizing can:

\begin{enumerate}
    \item To protect the fibre surface from damage
    \item To bind the fibres together in a yarn for ease of processing
    \item To impart anti-static properties
    \item To provide a chemical link between fibre surface and matrix
\end{enumerate}

\subsection{Carbon Fibre}

Carbon fibre is produced by the controlled \textbf{oxidation, carbonisation and graphitisation} of carbon-rich organic \textbf{precursors}(most common one is \textbf{PAN}) which are already in fibre form.

Carbon fibres can be grouped as:

\begin{enumerate}
    \item \textbf{High strength (HS)}
    \item \textbf{Intermediate modulus (IM)}
    \item \textbf{High modulus(HM)}
    \item \textbf{Ultra high modulus(UHM)}
\end{enumerate}

Carbon fibre has the highest \textbf{specific stiffness} of any commercially available fibre, in both tension and compression (in case of HS-fibre) and a high resistance to \textbf{corrosion, creep and fatigue}. 

Their impact strength, however, is lower than either glass or aramid, particularly HM and UHM fibres.

\fastfigure{PAN.png}

\subsubsection*{Carbon Fibre Production by PAN}

Tensile modulus of CF can mainly be determined by \textbf{carbonization/graphitization temperature}. Also, the nature of precursors (\textbf{PAN}, \textbf{pitch}) affects on the easiness of reaching a certain level modulus.

Some clothes says "acrylic", is a blended copolymer of \textbf{polyacrylonitrile}, with low purity and short chain.

Here we show the effect of heat treatment on strength and Young's modulus.

\fastfigure{Heat_vs_Young.png}

\subsubsection*{Carbon Fibre Production by Mesophase Pitch}

Heating to 
\begin{center}
    Single BSU $\to$ Distorted Columns $\to$ Wrinkled Layers $\to$ Stiff Layers
\end{center}

\fastfigure{pitch.png}

Graphitisation between 1500 and 3000$^\circ \text{C}$ gives improved orientation and properties. Heating above 2000$^\circ \text{C}$ gives higher \textbf{crystallinity} and higher (anisotropic) modulus - \textbf{HM-Carbon}.

\subsection{Aramid Fibres}

Aramid fibre is a man-made organic polymer produced by spinning from a \textbf{liquid crystalline solution}.

\begin{enumerate}
    \item High strength and low density.
    \item Good resistance to impact, abrasion, chemical and thermal degradation.
    \item Low elongation to break.
    \item Low electrical conductivity.
    \item \textbf{POOR} compressive strength, and degradate under UV light.
\end{enumerate}

Aramid fibres are highly \textbf{anisotropic - 1D} - with strong \textbf{covalent bonds} in the fibre direction and weak \textbf{hydrogen bonding} between the chains. 

This highly anisotropic character results in a \textbf{low shear modulus} and weak transversal and compressive properties.

\subsection{Polyethylene Fibres}

High-performance polyethylene fibres - like \textbf{Dyneema} or \textbf{Spectra} - are the world's strongest and lightest fibres. 

It has outstanding toughness, low density and high resistance to chemicals, water, and ultraviolet light, and can absorb extremely high amounts of energy.(To use these properties, a unidirectional sheet called \textbf{Dyneema UD} or \textbf{Spectra Shield} was developed)

It has excellent \textbf{vibration damping}, \textbf{flex fatigue} and internal fiber-friction characteristics, and low \textbf{dielectric constant} makes it virtually transparent to radar.

\fastfigure{polyethylene_produce.png}

\textbf{Solution (gel) spinning} process. The highly drawn fibre contains an almost $100\%$ \textbf{crystalline structure} with nearly perfectly arranged molecules, which promotes its extremely high strength, modulus, and other excellent properties.

\begin{mynote_env}[frametitle=Free breaking length]
    This is the theoretical length at which a fiber will break under its own weight.\\
    The free breaking length is independent of the size of the fiber and is determined only by the material properties. \textbf{HPPE} has higher free breaking length than polyester or polyamide.
\end{mynote_env}


\subsection{Boron Fibres}

Produced by \textbf{Chemical Vapour Depostion} on tungsten or carbon wire. Advantages of boron fibres are: 

\begin{enumerate}
    \item High modulus (400 GPa)
    \item High strength (3500 MPa)
    \item High temperature resistance (700 $^\circ \text{C}$)
    \item high compressive strength.
\end{enumerate}

However, it is so expensive that restrict its use.

\subsection{Ceramic Fibres}

Silicon carbide fibres are also made by \textbf{CVD} using a tungsten wire as substrate. However usually they are in the form of very short '\textbf{whiskers}' (\textbf{Nicalon}) and are mainly used in areas requiring high temperature resistance.

They are more frequently associated with \textbf{non-polymer matrices} such as metal alloys.

\subsection{Natural Fibres}

Natural fibres (flax, jute, hemp, sisal) are \textbf{renewable}, cheap (\texteuro0.35-1.5/kg), \textbf{lightweight}, and \textbf{biodegradable}.

\begin{mynote_env}[frametitle=Advantages and Disadvantages vs Glass]
    \textbf{Advantages}:

    Low density (1.3-1.5 g/cm$^3$), no skin irritation, non-abrasive to equipment, CO$_2$ neutral.\\
    \textbf{Disadvantages}: 
    
    Moisture sensitive, limited heat resistance, irreproducible quality, poor impact and compressive strength (fail by \textbf{kink-band formation}).
\end{mynote_env}


\begin{table}[h]
    \centering
    \begin{tabular}{lcccc}
        \hline
        Property & Flax & Jute & Sisal & Glass \\
        \hline
        Density (g/cm$^3$) & 1.54 & 1.46 & 1.33 & 2.55 \\
        Strength (MPa) & 600-1500 & 400-800 & 600-700 & 2400 \\
        Stiffness (GPa) & 40-70 & 10-30 & 38 & 73 \\
        Price (\texteuro/kg) & 0.5-1.5 & 0.35 & 0.6-0.7 & 1.3 \\
        \hline
    \end{tabular}
\end{table}

Applications include press moulding parts and wood fibre plastics ("mouldable wood").

\subsection{Textile Constructions}

\subsubsection*{Fibre Finishes}

To allow handling with minimum damage and to promote \textbf{fibre/matrix interfacial bond strength}.

\textbf{Glass Fibre Finishes}: 'Dual-function' finish for direct fibre processes (prepregging, pultrusion, filament winding). For weaving: protective finish (starch-based) applied first, then scoured and treated with matrix-compatible finish.

\textbf{Carbon Fibre Finishes}: Epoxy based, 1-2\% for weaving, 0.5-1\% for tape prepregging or filament winding.

\textbf{Aramid Fibre Finishes}: Composite finish, rubber compatible finish (belts and tyres), waterproof finish (ballistic soft armour).

\subsubsection*{Non-woven Mat}

\textbf{Chopped strand mat (CSM)} and \textbf{continuous strand mat} in non-woven, random pattern for equal distribution in all directions.

\begin{itemize}
    \item \textbf{Chopped strand mat}: Strands chopped to various lengths, deposited randomly. Used for hand lay-up, continuous laminating.
    \item \textbf{Continuous strand mat}: Deposited in swirl pattern. Used for press molding, RTM, pultrusion.
\end{itemize}

Available in weights from 200 to 1000 g/m$^2$.

\subsubsection*{Woven Fabrics}

\textbf{Woven roving}: Heavy, coarse, drapable fabric of plain weave. Alternated with CSM in multi-layer laminates to increase in-plane shear and peel strength.

\textbf{Weave styles}:
\begin{itemize}
    \item \textbf{Plain weave}: Each warp fibre passes alternately under and over each weft fibre. Symmetrical, good stability, most difficult to drape, high crimp lowers mechanical properties.
    \item \textbf{Twill weave}: One or more warp fibres weave over and under two or more weft fibres. Superior wet out and drape, smoother surface, slightly higher mechanical properties.
    \item \textbf{Satin weave}: Fewer intersections of warp and weft (harness number 4, 5, 8). Very flat, good wet out, high drape, low crimp gives good mechanical properties.
\end{itemize}

\subsubsection*{Hybrid Fabrics}

Fabric with more than one type of structural fibre.

\begin{itemize}
    \item \textbf{Carbon/Aramid}: High impact resistance and tensile strength of aramid + high compressive stiffness of carbon. Both low density but high cost.
    \item \textbf{Aramid/Glass}: Low density, high impact resistance of aramid + good compressive and tensile strength of glass, lower cost.
    \item \textbf{Carbon/Glass}: High tensile/compressive stiffness of carbon + reduced density, glass reduces cost.
\end{itemize}

\subsubsection*{Unidirectional Fabrics}

Majority of fibres runs in one direction only ($\sim$10\% fibre in other direction to hold primary fibres in place).

\textbf{UD prepreg}: True UD fabrics offer highest possible fibre properties. Fibres are straight and uncrimped. Only resin system holds fibres in place.

\subsubsection*{Multi-axial Stitched Fabrics}

One or more layers of long fibres held in place by PET stitching thread.

\textbf{Advantages}: Fibres always straight and non-crimped; fewer layers needed; better wet-out and permeability.

\textbf{Disadvantages}: PET stitching doesn't bond well to some resins; slow production process; high machinery and fabric cost.

\subsubsection*{Braided Fabrics}

Fibres interlaced in spiral nature to form tubular fabric. Diameter controlled by number of fibres, spiral angle (25$^\circ$ to 75$^\circ$), intersections per unit length, and fibre size.

Used for masts, antennae, drive shafts and other tubular structures requiring torsional strength.

\section{Processing of Polymer Matrix Composites}

\fastfigure{Designandfabirc.png}

\subsection{Thermalset Resins}

\textbf{Unsaturated polyester (UP)}, \textbf{vinyl ester (VE)}, \textbf{epoxy (EP)}, \textbf{polyurethane (PUR)}.

\begin{enumerate}
    \item Generally supplied as liquid
    \item \textbf{Cross-linked} (cured) by chemicals and heat
    \item Heat reduces the instantanious viscosity
    \item Heat increases the rate of cure
    \item Cure increases the viscosity over time
    \item Product is a \textbf{3D molecular network}, whereas thermoplastic is usually a 1D chain
\end{enumerate}

\fastfigure{Temperature_vs_Viscosity.png}

\subsubsection*{Polyester Resin}

In polyester chain, the addition \textbf{styrene} can:

\begin{enumerate}
    \item Easy to handle by reducing its viscosity.
    \item \textbf{Crosslinking} the chain, cure the resin from liquid to solid.
\end{enumerate}

Polyester resin has limited storage life because it will \textbf{"gel"}. For use in moulding, a polyester resin requires the addition of \textbf{catalyst}, \textbf{accelerators} or \textbf{inhibitors}.

\subsubsection*{Vinylester Resin}

Vinylester resin is similar in structure to polyester, but the \textbf{reactive sites only positioned at end of chains}, the whole chain can deform, so this make vinylester \textbf{more ductile}.

The vinylester has \textbf{fewer ester group}, means that vinylesters exhibit \textbf{better resistance to water} and many other chemicals than their polyester counterparts.

Idealised chemical structure of a typical vinylester:

\vspace{1em}
\ce{{\chemfig{C=C-C(=[2]O)-[@{left,.5} ::0]C-C(-[2]OH)-C-O-Ph-C(-[2]C)(-[6]C)-Ph-O-[@{right,.5} ::0]C-C(-[2]OH)-C-O-C(=[2]O)-C=C} \polymerdelim[delimiters = {[]}, height=7pt, h align=false]{left}{right}}}

\section{Labs}

\subsubsection*{Breakdown}

\begin{enumerate}
    \item Coursework = $40\%$ of the Module
    \item Lab 1: Wax Composites = $50\%$ of Coursework.
    \item Lab 2: Composite Layup = $50\%$ of Coursework.
\end{enumerate}

\subsection{Wax Composites}

Create composites using wax and fiber, choose fiber type, orientation and combination.

Testing the sample by \textbf{Bend, Impact, Toughness}.

\subsection{Composite Layup for Aircraft Panels}

Build composite panel for $25kg$ overhead stowage bin, using epoxy + fiberglass with \textbf{manual hand layup}.

Testing the sample by bending setup and calculate stiffness(\textbf{EI}, by bending angle).

$$
\theta = PL^2/16EI
$$

$$
EI = PL^2/16\theta
$$

$$
I=(1/12)\times b \times h^3
$$

\section{Processing of Metal and Ceramic Matrix Composites}

The matrix material largely determines the processing method.

\begin{table}[h]
    \centering
    \begin{tabular}{c|c}
        \hline
        Composites Matrix & Upper operation temperature\\
        \hline
        polymer & $300^\circ \text{C}$\\
        metal(A) & $600^\circ \text{C}$\\
        ceramic & $>1000^\circ \text{C}$\\
        \hline
    \end{tabular}
\end{table}

\subsection{Metal Matrix Composites}


\subsection{Ceramic Matrix Composites}

Ceramic matrix composite has 2nd high service temperature only lower than \textbf{Alumina/Mullite-based refractories}.

The interaction between matrix and reinforcing fibres provide higher toughness, the mechanism is \textbf{cracks deflection at the matrix-fibre interface}.

\subsection{Carbon-Carbon Composites}

The carbon-carbon composite fabrication by \textbf{liquid or gas phase infilltration}.

\section{Interface}

The properties change gradually, so we should not talk about an interface, but an \textbf{interface region or interphase}.

The weaker interface can absorb energy through \textbf{pull-out and debonding}, so stronger interface composite will be more brittle.

\begin{mynote_env}[frametitle=Example Question]
    A carbon fibre composite of $1\, cm^3$ with a fibre volume fraction of $60\%$ and carbon fibres of $8$ microns, contains $1.2\, million$ fibres, with a total interfacial area of $3000\, cm^2$.
    \begin{itemize}
        \item Fibre volume: $V_f = 0.6 \times 1 = 0.6\,\text{cm}^3$
        \item Fibre diameter: $d = 8\,\mu\text{m} = 8\times10^{-4}\,\text{cm}$
        \item Fibre cross-sectional area: $A_f = \pi(d/2)^2 = \pi(4\times10^{-4})^2 \approx 5.03\times10^{-7}\,\text{cm}^2$
        \item Total fibre length: $L_{\text{total}} = V_f / A_f \approx 1.19\times10^6\,\text{cm}$
        \item Length per fibre (assuming continuous fibres): $L \approx 1\,\text{cm}$ (i.e., fibres run the full length of the $1\,\text{cm}^3$ cube)
        \item Verify number of fibres: $N = V_f / (A_f \cdot L) \approx 0.6 / (5.03\times10^{-7} \times 1) \approx 1.19\times10^6 \approx 1.2$ million $\checkmark$
        \item Interfacial area per fibre: $\pi d L = \pi \times 8\times10^{-4} \times 1 \approx 2.51\times10^{-3}\,\text{cm}^2$
        \item Total interfacial area: $N \times \pi d L \approx 1.2\times10^6 \times 2.51\times10^{-3} \approx 3000\,\text{cm}^2$ $\checkmark$
    \end{itemize}
\end{mynote_env}

Composites with \textbf{weak interfaces} are often \textbf{weaker} than the composites with \textbf{strong interfaces}, but strong interfaces composites are often \textbf{more brittle} than weak interfaces composites.

\textbf{Flexible interphases} like rubber coatings can lead to \textbf{tougher} composites.

\subsection{Interface Modification for Polymeric Fibres}

\begin{itemize}
    \item \textbf{Acid Treatment}: Introduce \textbf{functional groups}(\ce{-COH},\ce{-COOH},\ce{C=O}) to imcrease the adhesion.
    \item  \textbf{Plasma Treatment}: Use plasma to create functional groups on the fibre surface, improving adhesion.
    \item \textbf{Cornona Treatment}: Similar to plasma treatment, but uses a \textbf{high-voltage electrical discharge} to modify the fibre surface in normal pressure conditions.
\end{itemize}

\subsection{Failure Modes}

\begin{itemize}
    \item \textbf{Interfacial Debonding}. load is applied \textbf{perpendicular} to the fibre axis. No fiber or matrix fracture, \textbf{ideal} interfacial failure mode.
    \item \textbf{Matrix Deformation or Cracking}. before interfacial failure happen, the matrix deforms or fractures around the fibre. shows \textbf{fibre-matrix adhesion is strong than the matrix itself}.
\end{itemize}

\subsubsection*{Pull-out Test}

Embedding a single fibre in a matrix and applying a tensile load to pull the fibre out of the matrix.

\begin{equation*}
    \sigma_f \pi r = 2 \pi r L \tau \rightarrow \tau = \sigma r / 2L
\end{equation*}

\subsubsection*{ILSS - Interlaminar Shear Strength}

Testing the adhesion between layers. Interlaminar shear stress at mid-plane is calculated by:

\begin{equation*}
    \tau_{ILSS} = \frac{3FS}{4ab}
\end{equation*}

\subsubsection*{Fracture Toughness}

For given applied stress and pre-exiting crack size, the \textbf{Energy Release Rate:}

\begin{equation*}
    G=\frac{\sigma^2 \pi c}{E}
\end{equation*}

The critical energy release rate $G_c$ is a measure of the material's resistance to crack propagation.

Stress required to cause spontaneous crack growth can be calculated by:

\begin{equation*}
    \sigma_{*} = \sqrt[2]{\frac{G_c E}{\pi c}}
\end{equation*}

Stress Intensity Factor:

\begin{equation*}
    K = \sigma \sqrt{\pi c}
\end{equation*}

In critical value of $G_c$, $K_c$ named \textbf{fracture toughness} can be calculated by:

\begin{equation*}
    K_c = \sqrt[2]{G_c E}
\end{equation*}

\begin{table}[h]
    \centering
    \begin{tabular}{lll}
        \hline
        Method & Advantage & Disadvantage \\
        \hline
        Toughened Matrix & Easy & Viscosity\\
        Interleaving & Localised, No effect on viscosity & Lower $V_f$,$T_g$\\
        Z-pinning & Localised Reinforcement & Complex\\
        \hline
    \end{tabular}
\end{table}

\begin{figure}[h]
    \centering
    \includegraphics[width=0.6\textwidth]{Method1.png}
    \caption{Method Figure 1}
\end{figure}

\begin{figure}[h]
    \centering
    \includegraphics[width=0.6\textwidth]{Method2.png}
    \caption{Method Figure 2}
\end{figure}

\newpage
\section{Final Student Checklist \& Answer}

\subsubsection*{1. Explain the role of fibre, matrix and interface in a composite}

\textbf{Fibre} — provides \textbf{strength and stiffness}. Fibres carry the majority of the load because they have far higher mechanical properties than the unreinforced matrix. The fibre type, volume fraction, orientation and length all determine how effectively the load is transferred.

\textbf{Matrix} — serves three key functions:
\begin{enumerate}
    \item \textbf{Separates the fibres}, preventing brittle fibres from contacting and abrading each other.
    \item \textbf{Transfers load} between fibres via shear stress at the fibre--matrix interface.
    \item \textbf{Protects the fibres} from environmental attack (moisture, chemicals, UV) and mechanical damage during handling.
\end{enumerate}

\textbf{Interface (or Interphase)} — the region where fibre and matrix meet. A good interface ensures efficient \textbf{stress transfer} from matrix to fibre. The interface controls failure mode: a \textit{weak} interface promotes energy absorption via debonding and pull-out (tougher), while a \textit{strong} interface gives higher strength but may lead to brittle fracture.

\subsubsection*{2. Compare glass, carbon and aramid fibres}

\begin{table}[h]
    \centering
    \begin{tabular}{lccc}
        \hline
        Property & \textbf{Glass (E-glass)} & \textbf{Carbon (HS)} & \textbf{Aramid (Kevlar)} \\
        \hline
        Cost & Low (\texteuro1.3/kg) & High & Medium--High \\
        Stiffness & $\sim$73 GPa & 200--400 GPa (highest) & 70--180 GPa (lower) \\
        Tensile Strength & $\sim$2.4 GPa & 3--7 GPa & $\sim$3.6 GPa \\
        Impact Resistance & Poor & Poor (especially HM/UHM) & Excellent \\
        Compressive Strength & Good & Good (HS type) & Poor \\
        Density & 2.55 g/cm$^3$ (heavy) & $\sim$1.8 g/cm$^3$ & $\sim$1.44 g/cm$^3$ (light) \\
        Fatigue/Creep & Moderate & Excellent resistance & Good \\
        UV Resistance & Good & Good & Degrades under UV \\
        \hline
    \end{tabular}
\end{table}

\textbf{Applications}:
\begin{itemize}
    \item \textbf{Glass}: Boat hulls, chemical tanks, wind turbine blades, inexpensive structural parts — where low cost matters more than ultimate stiffness.
    \item \textbf{Carbon}: Aerospace structures, Formula 1, high-end sporting goods (bikes, golf clubs), aircraft panels — where maximum specific stiffness is critical.
    \item \textbf{Aramid}: Bullet-proof vests, helmets, tyre reinforcement, ropes/cables, sporting goods — where impact resistance and low weight are paramount, but compressive loads are minimal.
\end{itemize}

\textbf{Hybrids} (e.g., carbon/aramid, carbon/glass) combine the strengths of each fibre type to balance cost, stiffness and impact performance.

\subsubsection*{3. Why vinyl ester, polyester and epoxy matrices are selected for different environments}

\textbf{Unsaturated Polyester (UP)}:
\begin{itemize}
    \item \textit{Advantages}: Lowest cost, easy to handle (styrene reduces viscosity), room-temperature cure possible, widely available.
    \item \textit{Disadvantages}: High shrinkage ($\sim$8\%), limited chemical resistance (ester groups hydrolyse), limited storage life (gels over time), lower mechanical properties.
    \item \textit{Use}: General-purpose, non-critical applications — boat building, automotive body panels, storage tanks for non-aggressive media.
\end{itemize}

\textbf{Vinyl Ester (VE)}:
\begin{itemize}
    \item \textit{Advantages}: Better chemical resistance (fewer ester groups than polyester), more ductile (reactive sites only at chain ends allow chain deformation), good water resistance, higher toughness.
    \item \textit{Disadvantages}: More expensive than polyester, higher viscosity.
    \item \textit{Use}: Chemical storage tanks, pipes, marine structures exposed to water, environments requiring combined chemical + mechanical performance.
\end{itemize}

\textbf{Epoxy (EP)}:
\begin{itemize}
    \item \textit{Advantages}: Best mechanical properties (highest strength and stiffness), lowest shrinkage ($<$2\%), excellent adhesion to fibres, good chemical and temperature resistance, very low creep.
    \item \textit{Disadvantages}: Highest cost, longer cure cycles, often requires elevated temperature cure, moisture sensitivity during processing.
    \item \textit{Use}: Aerospace structures, high-performance sporting goods, carbon fibre prepregs, any application demanding maximum performance.
\end{itemize}

\textbf{Selection Rule of Thumb}: Polyester for cost-driven projects, vinyl ester for chemical/water exposure, epoxy where maximum mechanical performance is non-negotiable.

\subsubsection*{4. Thermosets vs Thermoplastics}

\begin{table}[h]
    \centering
    \begin{tabular}{p{4.5cm}p{5cm}p{5cm}}
        \hline
        Feature & \textbf{Thermoset} & \textbf{Thermoplastic} \\
        \hline
        Molecular Structure & 3D cross-linked network (irreversible) & 1D linear/branched chains (no cross-links) \\
        Cure Mechanism & Chemical reaction + heat $\to$ permanent set & Melt on heating, solidify on cooling (reversible) \\
        Recyclability & \textbf{Not recyclable} — cannot be re-melted & \textbf{Recyclable} — can be re-melted and re-formed \\
        Processing & Liquid resin infused into fibres, then cured & Must be heated above $T_m$ or $T_g$, high viscosity $\to$ difficult fibre wet-out \\
        Viscosity & Low (easy fibre impregnation) & High (requires high pressure/temperature) \\
        Cycle Time & Long (minutes to hours for cure) & Short (seconds to minutes — cooling only) \\
        Toughness & Generally brittle & Generally tougher, higher impact resistance \\
        Temperature Resistance & Good — does not soften on reheating & Limited — softens above $T_g$ or $T_m$ \\
        Storage & Limited shelf life (may pre-gel) & Indefinite shelf life \\
        Examples & Epoxy, polyester, vinyl ester, phenolic & PEEK, PPS, polypropylene, nylon \\
        \hline
    \end{tabular}
\end{table}

Thermoplastics have \textbf{very high melt viscosity} (100--1000 Pa·s vs $<$1 Pa·s for thermoset liquids), making it difficult to impregnate fibre bundles completely. This demands specialised techniques (e.g., commingled fibres, powder impregnation, high-pressure autoclave) and limits part size and fibre volume fraction.

\subsubsection*{5. Why 0\textdegree, 90\textdegree\ and $\pm$45\textdegree\ plies are used in laminates}

\textbf{0\textdegree\ plies} — fibres aligned with the primary load direction. Provide maximum \textbf{axial stiffness and strength} ($E_1$, $\sigma_1$). Used where tensile/compressive loads dominate (e.g., wing spar caps, beams in bending).

\textbf{90\textdegree\ plies} — fibres perpendicular to primary load. Provide \textbf{transverse strength} and resist \textbf{Poisson's contraction} effects. Prevent longitudinal splitting. Also carry loads in the secondary direction when biaxial loading exists.

\textbf{$\pm$45\textdegree\ plies} — fibres at 45\textdegree\ to the load axis. Provide \textbf{shear stiffness and strength} ($G_{12}$). Essential for resisting \textbf{torsional loads}, and provide \textbf{in-plane shear resistance}. In tension, they also contribute to \textbf{Poisson's ratio control}.

\textbf{Practical combination}: A \textit{quasi-isotropic} layup [0/$\pm$45/90]$_s$ gives near-isotropic in-plane properties — standard for structures experiencing multi-directional or unknown loads. The proportions are tuned based on the load case: more 0\textdegree\ for tension/compression-dominated parts, more $\pm$45\textdegree\ for shear/torsion-dominated parts.

\subsubsection*{6. Why balanced and symmetric laminates reduce twisting and improve reliability}

\textbf{Balanced laminate}: For every $+\theta$ ply, there is an identical $-\theta$ ply (same material, same thickness, same position relative to mid-plane). This cancels \textbf{shear--extension coupling} — a tensile load will not induce shear deformation.

\textbf{Symmetric laminate}: The stacking sequence is mirrored about the mid-plane (e.g., [0/90/0] vs [0/90]$_s$). This eliminates \textbf{bend--extension coupling} — an in-plane load will not cause the laminate to bend or twist.

\textbf{Consequences of un-balanced or un-symmetric layups}:
\begin{itemize}
    \item Un-balanced: tensile loading causes unwanted shear deformation ("scissoring").
    \item Un-symmetric: thermal cooling after cure causes \textbf{warpage} (residual stresses bend the part); in-plane loads induce out-of-plane bending.
\end{itemize}

\textbf{Why it improves reliability}: Balanced and symmetric laminates behave predictably under load — no unexpected twisting, warping, or coupling effects. This simplifies design analysis and ensures consistent, repeatable performance in service.

\subsubsection*{7. Select suitable processes for different production scenarios}

\begin{table}[h]
    \centering
    \begin{tabular}{p{2.5cm}p{3.5cm}p{4cm}p{4cm}}
        \hline
        \textbf{Scenario} & \textbf{Recommended Process} & \textbf{Why} & \textbf{Examples} \\
        \hline
        Low-volume, large parts & Hand lay-up / spray-up & Low tooling cost, simple equipment, flexible & Boat hulls, prototypes, wind blades \\
        Medium-volume, panels & Vacuum bagging / RTM & Better consolidation than hand lay-up, moderate tooling & Automotive panels, aircraft fairings \\
        High-volume, small parts & Compression moulding (SMC/BMC) & Very fast cycle time ($<$5 min), automated & Automotive body panels, electrical enclosures \\
        Hollow parts & Filament winding & Fibres wound precisely onto a mandrel, high $V_f$ & Pipes, pressure vessels, drive shafts \\
        Constant cross-section & Pultrusion & Continuous process, very high output rate, excellent fibre alignment & Rods, beams, channels, bridge decks \\
        Panel-like (flat/slight curvature) & Press moulding / RTM / prepreg autoclave & Controlled thickness, good surfaces on both sides & Aircraft panels, satellite structures \\
        Complex 3D shapes, low volume & Hand lay-up + autoclave or RTM & Drapable fabrics conform to complex tooling & Formula 1 chassis, prosthetic limbs \\
        \hline
    \end{tabular}
\end{table}

\textbf{Process Selection Key Factors}:
\begin{enumerate}
    \item \textbf{Production volume}: Low $\to$ hand lay-up; High $\to$ compression moulding/pultrusion.
    \item \textbf{Part geometry}: Hollow $\to$ filament winding; constant section $\to$ pultrusion; complex $\to$ RTM/hand lay-up.
    \item \textbf{Performance requirements}: Maximum properties $\to$ prepreg + autoclave; moderate $\to$ wet lay-up/RTM.
    \item \textbf{Cost constraints}: Cheapest $\to$ hand lay-up/spray-up; moderate $\to$ RTM; expensive $\to$ autoclave prepreg.
\end{enumerate}

\subsubsection*{8. How voids, dry spots, poor cure and fibre misalignment reduce performance}

\textbf{Voids} (entrapped air or volatile bubbles):
\begin{itemize}
    \item Act as \textbf{stress concentrators} — local stress $\gg$ nominal stress, initiating cracks.
    \item Reduce \textbf{interlaminar shear strength (ILSS)} — a 1\% void increase can reduce ILSS by $\sim$7\%.
    \item Reduce \textbf{compressive and flexural strength} significantly.
    \item Increase moisture uptake, leading to environmental degradation.
\end{itemize}

\textbf{Dry spots} (regions where resin has not penetrated):
\begin{itemize}
    \item Fibres carry \textbf{zero load} in that region — no stress transfer from matrix to fibre.
    \item Creates a \textbf{discontinuity} in the load path, concentrating stress at the dry spot boundary.
    \item Essentially a pre-existing flaw — drastically reduces strength and fatigue life.
\end{itemize}

\textbf{Poor cure} (under-cured or over-cured resin):
\begin{itemize}
    \item Under-cure: Incomplete cross-linking $\to$ low stiffness, low $T_g$, poor chemical resistance, creep-prone.
    \item Over-cure: Thermal degradation, embrittlement, colour change, reduced toughness.
    \item Improper mix ratio (for two-part systems): Stoichiometry off $\to$ unreacted components act as plasticisers.
\end{itemize}

\textbf{Fibre misalignment} (fibres not in the intended direction):
\begin{itemize}
    \item A misalignment of just \textbf{5\textdegree} can reduce longitudinal modulus by $\sim$10\%.
    \item Off-axis fibres create \textbf{unwanted shear--extension coupling}.
    \item In UD laminates, waviness (in-plane or out-of-plane) dramatically reduces compressive strength because fibres buckle more easily.
\end{itemize}

\subsubsection*{9. Calculate or explain $E_1$, $E_2$ and $G_{12}$ using rule-of-mixtures}

\textbf{Longitudinal Modulus $E_1$} (fibre direction — isostrain, parallel model):
\begin{equation*}
    E_1 = E_f V_f + E_m V_m \quad (\text{Voigt bound — upper bound})
\end{equation*}
Fibre and matrix experience the \textbf{same strain} ($\varepsilon_{1f} = \varepsilon_{1m} = \varepsilon_1$). The stiffer fibre dominates. This is the \textbf{rule of mixtures (RoM)}.

\textbf{Transverse Modulus $E_2$} (perpendicular to fibres — isostress, series model):
\begin{equation*}
    \frac{1}{E_2} = \frac{V_f}{E_f} + \frac{V_m}{E_m} \quad \text{or} \quad E_2 = \frac{E_f E_m}{E_f V_m + E_m V_f} \quad (\text{Reuss bound — lower bound})
\end{equation*}
Fibre and matrix experience the \textbf{same stress} ($\sigma_{2f} = \sigma_{2m} = \sigma_2$). The compliant matrix dominates. This \textbf{inverse RoM} is a poor predictor because it assumes uniform stress, which is unrealistic.

\textbf{Better Transverse Prediction — Halpin--Tsai Equations}:
\begin{equation*}
    E_2 = E_m \frac{1 + \xi \eta V_f}{1 - \eta V_f}, \quad \eta = \frac{(E_f/E_m) - 1}{(E_f/E_m) + \xi}
\end{equation*}
$\xi$ is a curve-fitting parameter ($\sim$2 for circular fibres in square array). This semi-empirical model accounts for non-uniform stress distribution.

\textbf{In-plane Shear Modulus $G_{12}$}:
\begin{equation*}
    G_{12} = G_m \frac{1 + \xi \eta V_f}{1 - \eta V_f}, \quad \eta = \frac{(G_f/G_m) - 1}{(G_f/G_m) + \xi}
\end{equation*}
Same Halpin--Tsai form, using shear moduli. $\xi \approx 1$ for shear loading of circular fibres.

\textbf{In-plane Poisson's Ratio $\nu_{12}$}:
\begin{equation*}
    \nu_{12} = \nu_f V_f + \nu_m V_m \quad (\text{RoM — generally accurate})
\end{equation*}

\subsubsection*{10. Short-fibre stress transfer, critical fibre length and interface load transfer}
    \textbf{Stress Transfer Mechanism}:
    Load is applied to the matrix, which transfers it to the fibre via \textbf{interfacial shear stress $\tau$} along the fibre surface. The tensile stress in the fibre builds up from zero at the ends to a maximum at the centre, following the \textbf{shear-lag model}.

    \textbf{Critical Fibre Length $L_c$}:
    The minimum fibre length required for the fibre to reach its ultimate tensile strength $\sigma_f^*$ at its centre before interfacial failure:
    \begin{equation*}
        L_c = \frac{\sigma_f^* d}{2 \tau_y}
    \end{equation*}
    where $d$ = fibre diameter, $\tau_y$ = interfacial shear strength (or matrix shear yield strength).

    \begin{itemize}
        \item $L \gg L_c$: Fibre is \textbf{continuous} — can achieve full reinforcing potential, fails by fibre fracture.
        \item $L > L_c$: Fibre is \textbf{long enough} — reaches $\sigma_f^*$ in the central region, contributes effectively.
        \item $L < L_c$: Fibre is \textbf{too short} — never reaches $\sigma_f^*$, fails by \textbf{pull-out} rather than fracture. Reinforcement efficiency is low.
    \end{itemize}

    \textbf{Fibre Orientation Effect}:
    \begin{itemize}
        \item Fibres aligned with load: maximum reinforcement (efficiency factor $\eta_0 \approx 1$).
        \item Random in-plane orientation: $\eta_0 \approx 3/8$ (2D random) or $\eta_0 \approx 1/5$ (3D random).
        \item The modified RoM for short-fibre composites: $E_c = \eta_0 \eta_L E_f V_f + E_m V_m$, where $\eta_L$ is a length efficiency factor ($\eta_L < 1$ for $L < L_c$).
    \end{itemize}

    \textbf{Interface Load Transfer}:
    Load transfer occurs through \textbf{shear at the interface}. The interfacial shear strength depends on: (1) chemical bonding, (2) mechanical interlocking (surface roughness), (3) residual thermal stresses (compressive clamping from differential shrinkage during cooling), and (4) fibre surface treatments (sizings, oxidation).

\subsubsection*{11. Wetting, interfacial bonding, debonding, matrix cracking and fibre pull-out}

\textbf{Wetting}:
The liquid resin must \textbf{spread over and penetrate} the fibre surface. Good wetting requires the resin surface energy to be lower than the fibre surface energy ($\gamma_{\text{resin}} < \gamma_{\text{fibre}}$). Measured by \textbf{contact angle} $\theta$ — lower $\theta$ = better wetting. Poor wetting $\to$ voids and dry spots.

\textbf{Interfacial Bonding Mechanisms}:
\begin{enumerate}
    \item \textbf{Chemical bonding}: Covalent bonds between fibre surface groups (e.g., \ce{-OH}, \ce{-COOH} from oxidation) and matrix reactive groups. Primary mechanism in glass (via silane coupling agents) and carbon (via surface oxidation + epoxy sizing).
    \item \textbf{Mechanical interlocking}: Resin penetrates surface roughness/porosity of the fibre, creating a mechanical key.
    \item \textbf{Physico-chemical}: Van der Waals forces, hydrogen bonding, acid--base interactions.
    \item \textbf{Residual stress}: Differential thermal contraction during cool-down creates radial compressive stress at the interface, enhancing frictional load transfer.
\end{enumerate}

\textbf{Failure Sequence} (from least to most severe):
\begin{enumerate}
    \item \textbf{Interfacial Debonding}: The fibre--matrix bond fails. Load is applied \textit{perpendicular} to the fibre axis. No fibre or matrix fracture — this is the \textbf{ideal} interfacial failure mode because it absorbs energy without catastrophic failure.
    \item \textbf{Matrix Cracking}: Before or after debonding, the matrix deforms plastically and then cracks around the fibre. Indicates \textbf{fibre--matrix adhesion is stronger than the matrix itself}. Cracks can bridge between fibres, forming a macroscopic crack path.
    \item \textbf{Fibre Pull-out}: After debonding, fibres are pulled out of the matrix against frictional resistance. This is a major \textbf{energy-absorbing mechanism} — the work done against friction contributes significantly to fracture toughness. Long fibres with strong interfaces fracture; short fibres below $L_c$ pull out.
    \item \textbf{Fibre Fracture}: Fibre reaches its ultimate tensile strength and breaks. This is the terminal event in continuous-fibre composites with strong interfaces. Minimises energy absorption — leads to brittle, catastrophic failure.
\end{enumerate}

\textbf{Energy Absorption Hierarchy}: Pull-out $>$ Debonding $>$ Matrix cracking $>$ Fibre fracture. This is why \textit{moderately weak} interfaces often produce \textit{tougher} composites than very strong interfaces.

\subsubsection*{12. Evidence-based retain, repair or replace decisions}

\textbf{Inspection Techniques}:
\begin{itemize}
    \item \textbf{Visual inspection}: Surface damage, discolouration, impact marks, fibre exposure, resin yellowing (UV degradation).
    \item \textbf{Tap test / Coin tap}: Audible change in pitch indicates delamination or disbond — quick, low-tech, effective for near-surface defects.
    \item \textbf{Ultrasonic C-scan}: Maps internal delaminations, voids and thickness variations. Quantitative depth and size information.
    \item \textbf{X-ray / CT}: Detects internal cracks, foreign object damage, fibre fracture.
    \item \textbf{Thermography}: Active heating + IR camera reveals sub-surface defects (delaminations block heat flow).
    \item \textbf{Shearography / Laser interferometry}: Detects near-surface disbonds and delaminations via surface strain anomalies under slight loading.
\end{itemize}

\textbf{Decision Criteria}:
\begin{table}[h]
    \centering
    \begin{tabular}{p{2cm}p{4.5cm}p{4.5cm}p{4cm}}
        \hline
        \textbf{Finding} & \textbf{Retain} & \textbf{Repair} & \textbf{Replace} \\
        \hline
        Damage extent & Cosmetic only (paint scratch, surface gel-coat chip) & Localised, within repairable area ($<$(10--15)\% of part area) & Widespread or structural ($>$15\% area) \\
        Delamination & None or $<$10 mm diameter, not at critical location & 10--50 mm, accessible, mid-thickness & $>$50 mm, multiple layers, near edges or joints \\
        Fibre fracture & None & Few surface fibres only & Fibre breakage through thickness (visible back-face damage) \\
        Moisture ingress & None (dry condition) & Slight, can be dried out before repair & Severe, matrix hydrolysed, $T_g$ depressed \\
        Service criticality & Non-structural or secondary structure & Primary structure with documented repair scheme (SRM) & Safety-critical, no approved repair \\
        Cost & — & Repair cost $<$ 50\% of replacement & Repair cost $>$ 50\% of replacement, or downtime unacceptable \\
        \hline
    \end{tabular}
\end{table}

\textbf{General Principle}:
\begin{itemize}
    \item \textbf{Retain} if the defect does not grow under service loads and is below the damage tolerance limit.
    \item \textbf{Repair} (scarf patch, doubler patch, resin injection) if damage is localised, accessible, and an approved repair procedure exists. Monitor after repair.
    \item \textbf{Replace} if damage compromises structural integrity, exceeds repairable limits, or no validated repair method exists for that material system and load case.
\end{itemize}

\newpage{}

\section{Unidirectional Laminate Stiffness}

\textbf{Microstructure Equation}: \\ \quad $\left[ \text{High stiffness, strength} \right]$ fibre + $\left[ \text{Good shear strength and low density} \right]$ matrix \\ \qquad = \textbf{High performance composite}.

\subsection{Fibre Contrent and Packing}

\subsubsection*{Purpose of Matrix}

The matrix serves to \textbf{separate the fibers}, prevent cracking and friction between brittle fibers and \textbf{transfer loads} between them.

\subsubsection*{Definition of Composite Microstructure}

A material made from two or more constituent materials with significantly different physical or chemical properties that, when combined, produce a material with characteristics different from the individual components.

\subsubsection*{Multi-Level Mechanics of Composites}

\begin{equation*}
    E^* = f(E_f,E_m,V_f,S,A)
\end{equation*}

Where,(should be check, because it is auto-completed by AI)
\begin{itemize}
    \item $E^*$ is the effective stiffness of the composite
    \item $E_f$ is the stiffness of the fiber
    \item $E_m$ is the stiffness of the matrix
    \item $V_f$ is the volume fraction of the fiber
    \item $S$ is the fiber orientation and distribution
    \item $A$ is the interfacial adhesion between fiber and matrix
\end{itemize}

The density of the composite can be calculated by:
\begin{equation*}
    \rho_c = \varphi_f \rho_f + \varphi_m \rho_m
\end{equation*}

Where,
\begin{itemize}
    \item $\rho_c$ is the density of the composite
    \item $\varphi_f$ is the volume fraction of the fiber
    \item $\rho_f$ is the density of the fiber
    \item $\varphi_m$ is the volume fraction of the matrix
    \item $\rho_m$ is the density of the matrix
\end{itemize}

\subsection{Rules of Mixtures}

\textbf{Assumption}: Fibres are uniform, parallel, and perfectly bonded to the matrix. Longitudinal load produces.

\subsubsection*{Axial Stiffness}

Deformed in the axial or fibre direction, the fibre and matrix are \textbf{in parallel}, so they experience the \textbf{same strain}:

\begin{equation*}
    e_{1f} = e_{1m} = e_1
\end{equation*}

Inverse rule of mixtures for axial stiffness:
\begin{equation*}
    \frac{1}{E_c} = \frac{V_f}{E_f} + \frac{V_m}{E_m}
\end{equation*}

\subsection{Halpin-Tsai set of Equations for Transverse}

This \textbf{semi-empirical} model takes into account the \textbf{non-uniform stress distribution} in the transverse direction and provides prediction of the composite's transverse properties.

\begin{equation*}
    E_2 = \frac{E_m(1-\xi \varphi_f)}{1-\mu \varphi_f}
\end{equation*}

\begin{equation*}
    G_{12} = \frac{G_m (1+\xi \mu V_f)}{1-\mu V_f}
\end{equation*}

\subsection{Poisson's Ratio}

Most material will be thinner in the transverse direction when stretched in the axial direction, so the \textbf{Poisson's ratio} is defined as:

\begin{equation*}
    v = - \frac{\text{transverse strain}}{\text{axial strain}} = - \frac{e_2}{e_1}
\end{equation*}

\begin{equation*}
    v_{i,j} = \frac{e_j}{e_i}
\end{equation*}

\end{document}
